% ================================================================
% MortgageLoanDerivation.tex
% ================================================================
%
%  Used as an example of how to write a short mathematical article.

% ------------------------------------------ Initial Setup --------------------------------------

%  Run the "Typeset" command twice in TexShop, otherwise you won't get the table of contents.

\tracinglostchars=3  % Generate an error message if a UTF-8 character is not supported.


% ------------------------------ This document is an Article -------------------------------

\documentclass{article}

\author{ Sean E. O'Connor}
\title{Mortgage Loan Derivation}


% ------------------------ Packages for theorems et. al., tables, hyperlinks --------------------------------------
% For more packages, see  https://www.ctan.org/

\usepackage{amssymb}  % American Mathematical Society math symbols. https://www.ams.org/arc/resources/amslatex-about.html

\usepackage{amsthm}     % Help in writing theorems:  https://www.ams.org/arc/tex/amscls/amsthdoc.pdf

% Define styles for lemmas, theorems, definitions, remarks.

\newtheorem{lemma}{Lemma}[section]  % variable name, caption, section number.
\theoremstyle{definition} % Bold title, Roman text body.

\newtheorem{theorem}{Theorem}[section]
\theoremstyle{definition} % Bold title, Roman body.
\RenewDocumentCommand{ \qedsymbol }{}{ \(\blacksquare\)  }  % Redefine the Q.E.D. symbol as a mathematical black square.

\newtheorem{remark}{Remark}
\theoremstyle{remark} % Italic title, Roman body.

% Tables of good quality: https://ctan.mirrors.hoobly.com/macros/latex/contrib/booktabs/booktabs.pdf

\usepackage{booktabs}

% Hyperlinks:  https://ctan.math.illinois.edu/macros/latex/contrib/hyperref/doc/hyperref-doc.html

\usepackage[
  colorlinks
]{hyperref} % Hyperlinks.  This must be the last package loaded.

\hypersetup{   % Which colors to use.
   urlcolor=red,
   citecolor=purple,
   linkcolor=blue
 }
  
% ----------------- Document begins:  title, abstract, table of contents --------------------

\begin{document}

\maketitle

\begin{abstract}
We derive the equations for a \emph{mortgage} loan and give several useful formulas.
\end{abstract} 

\footnote{I'm using \LaTeX and \TeX Shop.  See \url{https://www.ctan.org}. See source code for more documentation.} 

\tableofcontents

\bigskip % Put extra line space after the contents.

% ------------------------------------------ First section begins --------------------------------------

\pagebreak 

\section{The Repayment Structure of a Mortgage Loan}

\subsection{Definition of Interest Rate}

The interest rate on the loan note is quoted as a yearly amount but is applied periodically for $m$ periods per year, giving

\begin{equation}
i =  {{Annual  \: interest    \:  percent / 100} \over m}
\end{equation}

Most house and car loans have a monthly period, $m = 12$.    % Inline math can use dollar sign delimiters like MathJax

\emph{NOTE:}  This is NOT the APR, which is the interest percentage as if the loan were compounded yearly instead of monthly.  I would ignore the APR.

\subsection{Periodic Payments of Principal + Interest}

The $j^{th}$ payment PMT consists of interest $I_j$ on the previous unpaid balance $BAL_{j-1}$

This payment is a fixed amount every month by design.

Only the remainder of the payment after interest  $R_j$ actually reduces the principal of the loan.

Payments from the initial payment to the end payment have the form,

\bigskip % Put extra line space before and after the table.

\begin{tabular}{ clll }
\toprule
Payment Number & Payment Amount  & Interest & Balance \\
\midrule
${1^{st}}$   &   $PMT = R_1 + I_1$   &   $I_1 = i \: BAL_0$   &  $BAL_1 = BAL_0 - R_1$  \\
$2^{nd}$    &   $PMT = R_2 + I_2$   &  $I_2 = i \: BAL_1 $   & $BAL_2 = BAL_1 - R_2$ \\
\ldots & \ldots & \ldots & \ldots \\
$j^{th}$      &  $PMT = R_j + I_j$      & $I_j = i \: BAL_{j-1}$  &  $BAL_j = BAL_{j-1} - R_j$ \\
\ldots & \ldots & \ldots & \ldots \\
$n^{th}$      &  $PMT = R_n + I_n$  & $I_n = i \: BAL_{n-1}$ &  $BAL_n= BAL_{n-1} - R_n = 0 $ \\
\bottomrule
\end{tabular}

\bigskip

n = the number of periods of the loan

\begin{remark}
The loan amount is $PV = BAL_0$ which is the original balance before any payments are made.
\end{remark}

\begin{remark}
The last balance is 0 because the loan is paid off, $BAL_n = 0$.
\end{remark}

We could write an Excel, Perl or Python script to compute this, but there are closed form solutions as we'll show.

\pagebreak 

\section{Closed Form Solution}

\subsection{Difference Equation}

Substituting we get a difference equation for the balances in terms of interest i and constant periodic payment PMT,

\begin{equation}
BAL_j = BAL_{j-1} - PMT + I_j
\end{equation}

\begin{equation}
BAL_j = BAL_{j-1} - PMT + i BAL_{j-1}
\end{equation}

\begin{equation}
BAL_j = (1+i) BAL_{j-1} - PMT
\end{equation}

Let's solve this difference equation.  But first, some lemmas,

\begin{lemma}{}
The geometric progression

\begin{equation}
g = 1 +  b + b^2 + \ldots + b^{n-1}
\end{equation}

has the formula,

\begin{equation}
g = {{1 - b^n} \over {1 - b}}
\end{equation}

\end{lemma}

\begin{proof}

Multiply by b and subtract,

\begin{equation}
g - b g = g(1 - b) = \left( 1 +  b + b^2 + \ldots + b^{n-1}  \right) - \left( b +  b^2 + b^3 + \ldots + b^n \right) = 1 - b^n
\end{equation}

to get

\begin{equation}
g = {{1 - b^n} \over {1 - b}}
\end{equation}

\end{proof}

\begin{lemma}{}
 If

\begin{equation}
u_k = a + b u_{k-1}
\end{equation}

then 

\begin{equation}
u_n = a {1 - b^n  \over 1-b } + b^n u_0
\end{equation}

\end{lemma}

\begin{proof}   Expand out terms to see the pattern and use mathematical induction if you like,

\begin{equation}
u_1 = a  + b u_0
\end{equation}

\begin{equation}
u_2 = a  + b u_1 = a + b(a + b u_0) = a (1 +  b) + b^2 u_0
\end{equation}

\begin{equation}
u_3 = a  + b u_2 = a + b [ (1+b) + b^2 u_0] = a (1 +  b + b^2) + b^3 u_0
\end{equation}

The general formula is

\begin{equation}
u_n = a (1 +  b + b^2 + \ldots + b^{n-1}) + b^n u_0
\end{equation}

Now use the geometric progression lemma above.

\end{proof}

\subsection{Closed Form Solution for $BAL_k$}

To solve the balance difference equation let $u_k = BAL_k$, $a = -PMT$, and $b = i+1$

The balance AFTER payment k is then

\begin{equation}
BAL_k = {(-PMT) {1 - (1+i)^k \over 1 - (1+i) } + (1+i)^k BAL_0}
\end{equation}

rearranging,

\begin{equation}
BAL_k =  {1 \over  (1+i)^{-k} }   \left(  PV + PMT  { (1+i)^{-k} -1 \over i }   \right)
\end{equation}

Now set $BAL_n = 0$ because the last nth payment PMT finishes the mortgage.
Check:  $BAL_0 = PV$

\pagebreak 

\section{Other Useful Closed Form Solutions}

\subsection{Periodic Payment PMT From PV, n and i}

The constant periodic payment is then

\begin{equation}
PMT =  {i PV \over  1-(1+i)^{-n}} 
\end{equation}

\subsection{Number of Periods n from PV, i and PMT}

The number of periods (months) comes from the loan amount PV and the periodic payment PMT after inverting,

\begin{equation}
n = -{ ln \left(     {1 - i {PV \over PMT}}    \right) \over { ln( 1 + i ) } }
\end{equation}

where we're using the natural logarithm, but any log will work.

\subsection{Reduced Number of Periods $n'$ from Increasing the PMT}

You can use this to determine the reduced number of periods if you prepay principal.  Change $PMT$ to $(PMT + additional\:  principal \: prepayment)$

\begin{equation}
n' = -{ ln \left(     {1 - i {PV \over {PMT + additional\:  principal \: prepayment}}}       \right) \over { ln( 1 + i ) } }
\end{equation}

\subsection{Equivalent Simple Interest $i'$}

Equivalent simple interest $i'$ on the loan is how much interest would have been had it been payable up front immediately.

We use this formula to deduce it,

\begin{equation}
PV + i' PV = n PMT
\end{equation}

Equivalent simple interest is then

\begin{equation}
i' =  {n i \over  1-(1+i)^{-n}} -1
\end{equation}

\subsection{Accumulated Interest Payments j though k}

What is the accumulated interest for payments j though k inclusive?

\begin{equation}
Int_{j \rightarrow k} = I_j + \ldots + I_k
\end{equation}

\begin{equation}
Int_{j \rightarrow k} = \left( k - j + 1 \right) PMT - \sum_{l=j}^k R_l
\end{equation}

but

\begin{equation}
BAL_l - BAL_{l-1} = -R_l
\end{equation}

We get a telescoping sum

\begin{equation}
- \sum_{l=j}^k R_l = BAL_k - BAL_{j-1}
\end{equation}

Interest AFTER the kth payment is

\begin{equation}
Int_{1 \rightarrow k}  = k PMT + BAL_k - PV
\end{equation}

Check:  When $k=n$,

\begin{equation}
Int_{1 \rightarrow n}  = n PMT + BAL_n - PV = n PMT - PV
\end{equation}

as expected because the total n payments of PMT include repaying PV plus interest.

\end{document}
